% !PRINTED 117
\weberpage{117}

We give $\sqrt R$ one of the two signs and, according to (5), obtain
\begin{equation}
\begin{array}{rclcrcl}
u^3&=&-\dfrac b2+\sqrt R,&& v^3&=&-\dfrac b2-\sqrt R,\\[0.75em]
u&=&\sqrt[3]{-\dfrac b2+\sqrt R},&& v&=&\sqrt[3]{-\dfrac b2-\sqrt R},
\end{array}
\tag{8}
\end{equation}
and therefore
\begin{equation}
y=\sqrt[3]{-\frac b2+\sqrt R}+\sqrt[3]{-\frac b2-\sqrt R}.
\tag{9}
\end{equation}

Multiplication of the two expressions (8) gives
\[
 u^3v^3=-\frac{a^3}{27}
\]
and shows that, if for $u$ a determinate one among the three values of the cube root is taken,
\begin{equation}
v=-\frac{a}{3u}=-\frac{a}{3\sqrt[3]{-\dfrac b2+\sqrt R}}
\tag{10}
\end{equation}
is in any case one of the three values of $\sqrt[3]{-\dfrac b2-\sqrt R}$. But it follows from (4) that only when we take this value for $v$ is equation (1) satisfied by $y=u+v$. Thus, corresponding to the three values of $u$, we obtain three roots of equation (1). From the existence of one root of the cubic equation, however, the existence of three already follows also from the fact that it has already been proved that the quadratic equation has two roots.

To derive the other two roots from the one root (9), put
\[
 \alpha=u+v
\]
and divide $y^3+ay+b$ by $y-\alpha$. The quotient whose roots are the two other roots $\beta,\gamma$ of the cubic equation (1) is
\[
 y^2+\alpha y+\alpha^2+a=0,
\]
whence, by (4) of the preceding paragraph, one finds
\[
 y=\frac{-\alpha\pm\sqrt{-3\alpha^2-4a}}{2}.
\]
If here one puts $\alpha=u+v$ and, by (4), $a=-3uv$, then it follows that
%CONT>>
